\section{Moving to Mathlib}\label{sec:13-next-steps:02-moving-to-mathlib:1}

Everything here was reinvented on purpose instead of imported, so that every moving part would be visible. Mathlib, Lean's community mathematics library, already has much more general and well-tested versions of all of this:

\begin{itemize}
\tightlist
\item
  \texttt{Mathlib.Algebra.Group.Defs} --- \href{https://loogle.lean-lang.org/?q=Group}{\texttt{Group}}, \texttt{CommGroup}, built with Lean's \textbf{type class} mechanism (\texttt{class ... extends ...}) instead of a plain \texttt{structure}. This lets notation like \texttt{a * b}, \texttt{a⁻¹}, \texttt{1} work the same way across every group, without threading a \texttt{Grp :} argument through every definition by hand, and lets Lean find instances automatically through typeclass search.
\item
  \texttt{Mathlib.Algebra.Ring.Defs} --- \href{https://loogle.lean-lang.org/?q=Ring}{\texttt{Ring}}, \texttt{CommRing}, \texttt{Field}, and the whole hierarchy in between (\texttt{Semiring}, \texttt{NonUnitalRing}, \ldots).
\item
  \texttt{Mathlib.Algebra.Module.Defs} --- \href{https://loogle.lean-lang.org/?q=Module}{\texttt{Module}}, much more general than Chapter 10's hand-built version, with the entire linear-algebra library (\texttt{Mathlib.LinearAlgebra.*}: bases, dimension, tensor products, exact sequences) built on top.
\item
  \texttt{Mathlib.Combinatorics.Quiver.Basic} and \texttt{Mathlib.Algebra.Category.*} --- quivers as the underlying data of a category (a category is ``a quiver plus identities and composition satisfying associativity,'' the same free-category construction from Chapter 11), and \texttt{Mathlib.CategoryTheory} more broadly.
\item
  Path algebras specifically show up in representation-theory-oriented corners of Mathlib and in dedicated Lean projects on quiver representations. Searching Mathlib's docs for ``quiver'' and ``path'' is a good starting point once the type-class style is familiar.
\end{itemize}

The jump from this book's \texttt{structure}-based definitions to Mathlib's \texttt{class}-based ones is mostly about \textbf{ergonomics}: automatic instance resolution, shared notation, and inheritance diamonds (the ambiguity that arises when a structure extends two parents with a common ancestor) already resolved. It is not really about mathematical content --- the axioms already learned are the same axioms, merely packaged so Lean's elaborator can find them without a \texttt{Grp} argument named in every theorem.

\subsection{Two theorems for free}\label{sec:13-next-steps:02-moving-to-mathlib:2}

Everything above is a promise that Mathlib is \emph{more general}. Here are two concrete payoffs --- things this book explicitly could not state, using exactly the examples already built.

\textbf{\texttt{ZMod 3} is a field, not just a ring.} Chapter 8 §5 built \texttt{fin3Ring} and said in so many words that a \texttt{Field} would need every nonzero element to be invertible, ``true for \(\mathbb{Z}/3\) precisely because \(3\) is prime, but not part of \texttt{Ring}'s axioms and not checked here.'' Mathlib already has this instance, for real:

\begin{lstlisting}
example : Field (ZMod 3) := inferInstance

-- A fact `fin3Ring` (a mere `Ring`) genuinely cannot state: every
-- nonzero element has a multiplicative inverse.
example : (*$\forall$*) a : ZMod 3, a (*$\neq$*) 0 (*$\rightarrow$*) (*$\exists$*) b, a * b = 1 := by decide
\end{lstlisting}

\textbf{Lagrange's theorem, applied to the non-abelian example from Chapters 6-7.} \texttt{Equiv.Perm (Fin 3)} (Chapter 6 §4's Mathlib analogue of \texttt{perm3Group}) never had subgroups defined for it --- the book built one group, not the lattice of its subgroups. Mathlib's \texttt{Subgroup} type already comes with Lagrange's theorem attached, so applying it to a real subgroup costs nothing beyond naming the subgroup:

\begin{lstlisting}
-- Lagrange's theorem, fully generic: a subgroup's size divides the
-- ambient group's size.
example (s : Subgroup (Equiv.Perm (Fin 3))) :
    Nat.card s (*$\mid$*) Nat.card (Equiv.Perm (Fin 3)) :=
  Subgroup.card_subgroup_dvd_card s

-- Concretely: |S_3| = 3! = 6...
example : Nat.card (Equiv.Perm (Fin 3)) = 6 := by
  rw [Nat.card_eq_fintype_card, Fintype.card_perm, Fintype.card_fin]
  rfl

-- ...and the cyclic subgroup generated by the 3-cycle finRotate 3 has
-- order 3, since (finRotate 3)^3 = 1 and finRotate 3 (*$\neq$*) 1 with 3 prime.
example : Nat.card (Subgroup.zpowers (finRotate 3)) = 3 := by
  rw [Nat.card_zpowers]
  apply orderOf_eq_prime
  (*$\cdot$*) decide
  (*$\cdot$*) decide
\end{lstlisting}

Thus \(3 \mid 6\) --- not asserted, but \emph{derived}, by instantiating a theorem Mathlib already proved once, generically, for every group. Neither of these two facts required a single new definition: both reuse objects this book already built, and both are facts this book's own from-scratch \texttt{Ring}/\texttt{Group} genuinely could not have stated, since it never built the surrounding machinery (invertibility, subgroups) that Mathlib already has. This is the concrete shape of ``moving to Mathlib'': not different mathematics, merely a much larger stock of already-proved consequences to draw on.
